% 小熊座（综述）
% license CCBYSA3
% type Wiki

本文根据 CC-BY-SA 协议转载翻译自维基百科\href{https://en.wikipedia.org/wiki/Ursa_Minor}{相关文章}。

小熊座（Ursa Minor，拉丁语意为“较小的熊”，与大熊座 Ursa Major 相对），又称小熊星座，是位于遥远北天的一组星座。与大熊座类似，小熊座的尾部也可被视作一只勺子的把手，因此在北美被称为“小北斗”（Little Dipper）：由七颗恒星组成，其中四颗构成勺状的斗部，形态与其对应的大北斗相似。小熊座是公元$^{[2]}$世纪天文学家托勒密列出的$^{[48]}$个星座之一，至今仍是现代$^{[88]}$个星座之一。由于北极星位于其中，小熊座在传统航海导航中，尤其对水手而言，具有重要意义。

北极星（Polaris）是该星座中最亮的恒星，为一颗黄白色超巨星，同时也是夜空中最亮的造父变星，其视星等在$^{[1.97]}$至$^{[2.00]}$之间变化。小熊座$\beta$星（Beta Ursae Minoris），又名 Kochab，是一颗已进入晚期演化阶段的恒星，因膨胀并冷却而成为一颗橙色巨星，视星等为$^{[2.08]}$，仅略暗于北极星。Kochab 与视星等约为$^{[3]}$等的小熊座$\gamma$星（Gamma Ursae Minoris）常被称为“北极星的守护者”或“极星守卫者”$^{[4]}$。在该星座中已有行星被探测到环绕其中四颗恒星运行，包括 Kochab。此外，小熊座还包含一颗孤立的中子星——Calvera，以及迄今发现的第二高温白矮星 H1504+65，其表面温度约为$^{[200{,}000]}$ K。
\subsection{历史与神话}
\begin{figure}[ht]
\centering
\includegraphics[width=6cm]{./figures/9a220c9be61db089.png}
\caption{天龙座环绕着小熊座，如约 1825 年在伦敦出版的一组星座图《乌拉尼娅之镜》（Urania's Mirror）中所描绘的那样。} \label{fig_Ursa_1}
\end{figure}